\section{서론}

\subsection{문제의식}
리듬 게임의 판정선 깜빡임, 경매 앱의 마감 직전 테두리 강조, 퀴즈 앱의 ``준비'' 애니메이션처럼, 인터페이스는 중요한 순간 직전에 시각 신호를 자주 보낸다.
이런 효과는 사용자가 더 빨리 반응하게 돕는다고 기대하고 쓰지만, 그 기대가 실제 사용 환경에서 맞는지는 대개 확인되지 않는다.
이 연구는 이 질문을 검증 가능한 한 문장으로 좁혀, 직접 만든 웹 실험으로 데이터를 모아 확인한다.

\subsection{이미 알려진 것}
주의 연구는 주의를 하나의 기능이 아니라 여러 하위 체계로 본다.
\citet{posner1990}은 경고 신호에 반응해 준비 상태를 높이는 \emph{각성}(alerting)을 정향, 실행 통제와 구분되는 기능으로 제안했고,
\citet{petersen2012}은 이후 20년의 연구를 바탕으로 이 구분을 갱신했다.
중립적 경고 신호, 즉 목표가 어디에 나올지나 무엇인지는 알려 주지 않는 신호 뒤에는 신호가 없을 때보다 반응이 빠르다 \citep{weinbach2012}.
\citet{fan2002}은 경고 단서가 있는 시행과 없는 시행의 반응시간 차이로 각성 효과를 재는 과제를 만들었다.
이 과제는 성인 40명에게서 개인 단위로도 신뢰할 만한 추정치를 냈다.
경고와 목표 사이 간격(예비 기간)이 단순 반응시간에 영향을 준다는 사실도 오래전부터 정리되어 있다 \citep{niemi1981}.
이 때문에 이 연구는 두 조건의 대기시간 분포를 똑같이 맞췄다(\ref{sec:design}절).

\subsection{웹으로 반응시간을 재도 되는가}
반응시간 연구는 전통적으로 실험실의 전용 장비에서 이뤄졌다.
전용 측정 장비로 웹 실험을 점검한 연구들은, 브라우저에서 측정한 반응시간이 실험실보다 변동이 조금 크지만
참가자 내 조건 비교에는 충분히 쓸 수 있다고 보고했다 \citep{reimers2015,bridges2020,anwylirvine2021}.
특히 \citet{pronk2020}은 터치스크린과 키보드 기기에서 반응시간이 일관되게 과대 측정되며 그 크기가 기기마다 다르다고 보고했다.
같은 연구는 이런 기기 간 차이가 절대 반응시간보다 같은 사람 안의 상대 반응시간에 덜 영향을 준다고 정리했다.
이 결과가 이 연구가 참가자 간 비교가 아니라 \emph{참가자 내} 설계를 택한 핵심 근거다.

\subsection{이 연구가 새로 확인하려는 것}
기존 연구가 보인 것은 (1) 통제된 실험실 환경에서 경고 신호가 반응을 빠르게 한다는 것과
(2) 웹 브라우저가 조건 간 반응시간 차이를 잡아낼 만큼 정밀하다는 것이다.
이 연구는 두 결과를 합쳐, 참가자가 \emph{자기 기기로, 감독 없는 환경에서} 하는 3분짜리 웹 과제에서도,
UI 효과처럼 생긴 \emph{화면 테두리 점등}이 같은 사람의 반응시간을 줄이는지 확인한다.
새로운 이론을 세우는 연구가 아니라, 알려진 효과를 현실적인 사용 환경에서 다시 확인하는 개념적 재현이다.

\subsection{가설}
\begin{quote}
\textbf{가설.} 목표 자극이 나타나기 500\,ms 전에 화면 테두리가 100\,ms 동안 밝아지는 시각 경고 신호가 있으면,
신호가 없을 때보다 같은 참가자의 단순 반응시간이 짧아진다.
\end{quote}
원인은 경고 신호의 유무이고, 결과는 단순 반응시간(ms)이며, 방향은 ``짧아진다''이다.
가설이 맞다면 대다수 참가자에게서 [신호 있음 중앙값 $-$ 신호 없음 중앙값]이 음수로 나오고,
사전 등록한 검정에서 그 중앙값이 0보다 작다는 결과가 나와야 한다.
가설이 틀렸다면 이 차이가 0 근처에 흩어지거나, 경고가 오히려 방해가 되어 양수가 많이 나온다.
